\documentclass[11pt, a4paper]{article}

% ── Packages ──────────────────────────────────────────────────────────
\usepackage[utf8]{inputenc}
\usepackage[T1]{fontenc}
\usepackage{lmodern}
\usepackage[margin=1in]{geometry}
\usepackage{amsmath, amssymb, amsthm}
\usepackage{booktabs}
\usepackage{longtable}
\usepackage{graphicx}
\usepackage{hyperref}
\usepackage{xcolor}
\usepackage{listings}
\usepackage{tcolorbox}
\usepackage{polya}

\hypersetup{
  colorlinks=true,
  linkcolor=polyablue,
  urlcolor=polyablue,
  citecolor=polyablue,
}

% ── Code listing style ────────────────────────────────────────────────
\lstdefinestyle{polya-python}{
  language=Python,
  basicstyle=\ttfamily\small,
  keywordstyle=\color{polyablue}\bfseries,
  commentstyle=\color{polyagray}\itshape,
  stringstyle=\color{polyagreen},
  numbers=left,
  numberstyle=\tiny\color{polyagray},
  numbersep=8pt,
  frame=single,
  framerule=0.4pt,
  rulecolor=\color{polyagray},
  breaklines=true,
  showstringspaces=false,
  tabsize=4,
}
\lstset{style=polya-python}
\title{Study Schedule}
\author{uber-polya}
\date{2026-02-26}

\begin{document}

\maketitle
\tableofcontents
\newpage

% ── Phase A: Problem Formulation ─────────────────────────────────────
\section{Problem Formulation}

\begin{polyamodel}

\textbf{Problem Type}: Find \hfill
\textbf{Domain}: Graph Theory / Combinatorics


\end{polyamodel}

% ── Universe ──────────────────────────────────────────────────────────
\subsection{Universe}

\begin{itemize}
  \item $Assignment: \{Linear Algebra, Abstract Algebra, Probability, Statistics, Real Analysis, Discrete Math\}$
  \item $Color Map: \{Linear Algebra, Abstract Algebra, Probability, Statistics, Real Analysis, Discrete Math\}$
  \item $Graph Stats: \{vertices, edges, max_degree, clique_number\}$
\end{itemize}

% ── Variables ─────────────────────────────────────────────────────────
\subsection{Variables}

\begin{longtable}{lll}
\toprule
\textbf{Symbol} & \textbf{Meaning} & \textbf{Type / Range} \\
\midrule
\endhead
$x_{ij}$ & 1 if item i assigned to slot j, 0 otherwise & $\{0, 1\}$ \\
\bottomrule
\end{longtable}

% ── Structure ─────────────────────────────────────────────────────────
\subsection{Structure}

A student has 6 subjects to study across 5 time slots per day (morning, late-morning, afternoon, late-afternoon, evening). Some subjects share prerequisites or overlapping content and should not be studied in adjacent time slots to avoid confusion. Create a weekly study timetable by assigning each subject to a time slot such that conflicting subjects never occupy neighboring slots.

% ── Mapping ───────────────────────────────────────────────────────────
\subsection{Real-World Mapping}

\begin{longtable}{ll}
\toprule
\textbf{Real-World Concept} & \textbf{Mathematical Object} \\
\midrule
\endhead
assignment & $x: solution vector (assignment)$ \\
total score/cost & $objective function value$ \\
\bottomrule
\end{longtable}

% ── Constraints ───────────────────────────────────────────────────────
\subsection{Constraints}

\begin{enumerate}
  \item $x_{ij} \in \{0, 1\} \quad \forall\, i, j$
  \hfill \textit{(binary decision variables)}
\end{enumerate}

% ── Objective / Claim ─────────────────────────────────────────────────
\subsection{Objective}

\begin{equation}
\text{optimize} \; f(x)
\end{equation}

% ── Approach ──────────────────────────────────────────────────────────
\subsection{Approach}

\begin{description}
  \item[Suggested approach:] Backtracking Distance-2 Graph Coloring
  \item[Complexity class:] 
  \item[Available tools:] Backtracking Distance-2 Graph Coloring
\end{description}
% ── Phase B: Solution ────────────────────────────────────────────────
\section{Solution}

\begin{polyaresult}

\textbf{Answer}: 6-element assignment: Linear Algebra -> Morning (8-10), Abstract Algebra -> Afternoon (1-3), Probability -> Morning (8-10), Statistics -> Afternoon (1-3), Real Analysis -> Afternoon (1-3), Discrete Math -> Morning (8-10)

\textbf{Objective Value}: $3$

\textbf{Optimal}: Yes \hfill
\textbf{Feasible}: Yes
\textbf{Algorithm}: Backtracking Distance-2 Graph Coloring \quad $$

\textbf{Time}: 2.6030989829450846e-05s

\textbf{Certificate}: Solver reports optimal.

\end{polyaresult}

% ── Solution Details ──────────────────────────────────────────────────
\subsection{Solution Details}

Assignment:
  Linear Algebra -> Morning (8-10)
  Abstract Algebra -> Afternoon (1-3)
  Probability -> Morning (8-10)
  Statistics -> Afternoon (1-3)
  Real Analysis -> Afternoon (1-3)
  Discrete Math -> Morning (8-10)

Objective value: z* = 3

Method: Graph coloring on a conflict graph. Vertices represent subjects, edges represent conflicts (subjects that should not be in adjacent slots). Colors represent time slots. Uses NetworkX's greedy\_color with the largest-first strategy. An adjacency constraint is enforced by expanding the conflict graph: if subjects A and B conflict, we add edges not just between A-B but also between any subjects whose assigned slots would be adjacent.

% ── Verification ──────────────────────────────────────────────────────
\subsection{Verification}

\begin{longtable}{lcl}
\toprule
\textbf{Check} & \textbf{Status} & \textbf{Detail} \\
\midrule
\endhead
Feasibility & \passmark & All constraints satisfied \\
Optimality & \passmark & Backtracking Distance-2 Graph Coloring \\
\bottomrule
\end{longtable}

% ── Phase C: Interpretation ──────────────────────────────────────────
\section{Interpretation}

% ── The Question & The Answer ─────────────────────────────────────────
\subsection{The Question}
A student has 6 subjects to study across 5 time slots per day (morning, late-morning, afternoon, late-afternoon, evening).

\subsection{The Answer}

\begin{polyainsight}[title={Bottom Line}]
The optimal solution has objective value z* = 3.
\end{polyainsight}

% ── What This Means ───────────────────────────────────────────────────
\subsection{What This Means}

- All 6 subjects assigned to time slots (colors 0-4)
- No two conflicting subjects share adjacent time slots
- A printable weekly timetable grid

Graph coloring on a conflict graph. Vertices represent subjects, edges represent conflicts (subjects that should not be in adjacent slots). Colors represent time slots. Uses NetworkX's greedy\_color with the largest-first strategy. An adjacency constraint is enforced by expanding the conflict graph: if subjects A and B conflict, we add edges not just between A-B but also between any subjects whose assigned slots would be adjacent.

% ── Sensitivity Analysis ──────────────────────────────────────────────

% ── Figures ───────────────────────────────────────────────────────────

% ── Recommendations ───────────────────────────────────────────────────
\subsection{Recommendations}

\begin{enumerate}
  \item The solution is provably optimal; no further improvement is possible within this model.
\end{enumerate}

% ── Limitations ───────────────────────────────────────────────────────
\subsection{Limitations}

\begin{polyawarnbox}
\begin{itemize}
  \item Model assumes deterministic parameters; real-world variability not captured.
  \item Solution quality depends on the accuracy of input data.
\end{itemize}
\end{polyawarnbox}

% ── Appendix ─────────────────────────────────────────────────────────

\end{document}